% AI-produced by GPT-6 Astra (OpenAI, ChatGPT Pro mode). No priority claim.
% This section uses only standard theorem/proof environments and explicit operators.
\subsection{Finite tests for totality: Lead 6}

\begin{lemma}[Known, with the reduction proved: finite-poset representatives]\label{FT:sphere-representatives}
Let $C$ be a small category and let $n\geq1$. Every free homotopy class of
maps $S^n\to |NC|$ has a representative $|Nv|:|ND|\to |NC|$, after a
homeomorphism $|ND|\cong S^n$, where $D$ is the poset of nonempty faces of
a finite triangulation of $S^n$ and $v:D\to C$ is a functor.
\end{lemma}
\begin{proof}
We recall the simplicial approximation argument, including the change of
domain that is needed here. Choose a finite ordered simplicial complex $K$
triangulating $S^n$. The natural weak equivalence
$NC\longrightarrow\operatorname{Ex}^{\infty}NC$ has Kan target.
Consequently a homotopy class $|K|\to|NC|$ is represented by a simplicial
map $f:K\to\operatorname{Ex}^{\infty}NC$. Since $K$ has finitely many
nondegenerate simplices, $f$ factors through
$\operatorname{Ex}^rNC$ for some finite $r$. The subdivision--extension
adjunction gives a simplicial map
$g:\operatorname{Sd}^rK\to NC$. Under realization the adjunction compares
$|g|$ with the original class by the last-vertex maps
$\operatorname{Sd}^rK\to K$; their realizations are homotopic to the
subdivision homeomorphisms. We may increase $r$ to ensure $r\geq1$.
Now $\operatorname{Sd}^rK$ is the nerve of the poset $D$ of nonempty faces
of the finite complex $\operatorname{Sd}^{r-1}K$. The nerve is fully
faithful on categories: the images of vertices and edges of $g$ specify
objects and arrows, and its two-simplices impose precisely the composition
equalities. Thus $g=Nv$ for a unique functor $v:D\to C$.
The facts about $\operatorname{Ex}^{\infty}$ and last vertices used in
this argument are \cite[Theorem~I.11.4, Section~III.4, and
Theorem~III.4.8]{GoerssJardine2009}; they concern arbitrary simplicial sets
and do not require $NC$ to be a simplicial complex.
\end{proof}

\begin{theorem}[Proved: Lead 6]\label{FT:finite-test}
Let $A$ be a small nonempty connected category. For a functor
$u:D\to A$, define the covariant $A$-set
\[
 H_u(a)=\pi_0(u\downarrow a),\qquad
 H_u(b)([d,f])=[d,bf]
\]
for $b:a\to a'$. Here $(d,f)$ has $f:u(d)\to a$, and an arrow
$(d,f)\to(d',f')$ is $\delta:d\to d'$ with $f'u(\delta)=f$.
Define a functor
\[
 j_u:D\longrightarrow\operatorname{El}_A(H_u),\qquad
 d\longmapsto(u(d),[d,1_{u(d)}]),
\]
whose value on an arrow is its image under $u$. The following conditions
are equivalent:
\begin{enumerate}
\item Every component of every $\operatorname{El}_A(H)$, for every
covariant $H:A\to\mathsf{Set}$, is weakly contractible.
\item For every finite poset $D$ and every $u:D\to A$, $Bj_u$ is
null-homotopic on every connected component of $BD$.
\item The same conclusion holds when $D$ is the face poset of a finite
triangulation of $S^n$, for every integer $n\geq1$.
\end{enumerate}
In particular the tests in (3) suffice; no bound on the dimension of the
spheres is being imposed.
\end{theorem}
\begin{proof}
The arrow $u(\delta):u(d)\to u(d')$ defines an arrow in the displayed
category of elements because $(d,u(\delta))$ and $(d',1)$ are joined by
$\delta$ in $u\downarrow u(d')$. This proves that $j_u$ is well-defined.

The essential factorization is strict, not merely a homotopy
factorization. Let $H:A\to\mathsf{Set}$ and let
$v:D\to\operatorname{El}_A(H)$ be a functor. Write
$v(d)=(u(d),h_d)$, where $u=\pi v$ and
$\pi:\operatorname{El}_A(H)\to A$ is the projection. Define
\[
 \eta_v:H_u\longrightarrow H,\qquad
 (\eta_v)_a([d,f])=H(f)(h_d).
\]
If $f'u(\delta)=f$, then
$H(f)(h_d)=H(f')H(u(\delta))(h_d)=H(f')(h_{d'})$, so this formula is
constant on components of $u\downarrow a$. It is natural in $a$ because
$H(b)H(f)(h_d)=H(bf)(h_d)$. It therefore induces a functor over $A$,
\[
 E_v:\operatorname{El}_A(H_u)\longrightarrow
       \operatorname{El}_A(H),\qquad
 (a,\xi)\longmapsto(a,(\eta_v)_a(\xi)),
\]
which is the identity on underlying arrows of $A$. Direct substitution
shows $E_vj_u=v$ as functors.

Assume (1). Each component of the target of $Bj_u$ is a weakly
contractible CW complex and hence contractible by Whitehead's theorem.
The continuous image of a component of $BD$ lies in one such component,
so the restricted map is null-homotopic. This proves (2), and (2)
implies (3).

Assume (3), fix $H$ and a component $C$ of $\operatorname{El}_A(H)$,
and consider any map $f:S^n\to BC$, $n\geq1$. By
Lemma~\ref{FT:sphere-representatives}, after identifying the source with
$BD$ its free homotopy class is represented by $Bv$, for a functor
$v:D\to C$. The preceding factorization gives $Bv=BE_v\,Bj_{\pi v}$.
Condition (3) makes this composite null-homotopic. Thus every unbased
sphere map into the connected space $BC$ is null. This implies the
vanishing of every based homotopy group: for $n=1$ a free null-homotopy
makes a loop conjugate to the identity and hence trivial; for $n>1$ the
change of basepoint acts by automorphisms and fixes the zero class.
Therefore $BC$ is weakly contractible. The case $H=\varnothing$ has no
components and imposes no extra condition. Empty or disconnected $D$
causes no difficulty in (2), which is expressly componentwise.
\end{proof}

\subsection{The opposite interval-type category: Lead 8}

\begin{proposition}[Proved: precise variance and a homological obstruction]\label{FT:ty}
Let $\mathsf{Ty}$ be the category in PRETOT Definition E.1: its objects
are finite words $\tau$ in $+,-$, including the empty word, and a morphism
$\tau\to\rho$ is an endpoint-preserving functor
$I_\tau\to I_\rho$ whose function on vertices is nondecreasing. Here
$I_\tau$ is the free category on the oriented line prescribed by $\tau$.
Then $\mathsf{Ty}^{\mathrm{op}}$ is isomorphic to the category of these
words and nondecreasing sign-preserving maps of their letter sets.

Put $p=(++)$, $q=(+-+)$, and use one-based letter indices. The arrows
\[
 e=(1,3),\quad c=(1,1):p\rightrightarrows q
\]
have no completion to a commutative square. For every nonzero unital
commutative ring $R$ they have no affine equalizer in
$\operatorname{Aff}_R(\mathsf{Ty}^{\mathrm{op}})$.
In fact, for the covariant set quotient
\[
 H=\mathsf{Ty}^{\mathrm{op}}(p,-)/(ge\sim gc)
\]
we have
$H_1(B\operatorname{El}_{\mathsf{Ty}^{\mathrm{op}}}(H);\mathbb Z)
\cong\mathbb Z$. Thus $\mathsf{Ty}^{\mathrm{op}}$ is neither total nor
exact. The phrase ``over any ring'' in Lead 8 must exclude the zero ring;
with that qualification the lead is true.
\end{proposition}
\begin{proof}
For $\tau$ of length $n$ and $\rho$ of length $m$, write a type map on
vertices as
$\phi:\{0,\ldots,n\}\to\{0,\ldots,m\}$, with
$\phi(0)=0$, $\phi(n)=m$ and $\phi(i-1)\leq\phi(i)$. The functor
condition says exactly that every letter of $\rho$ in the interval
$(\phi(i-1),\phi(i)]$ has the sign of the $i$th letter of $\tau$;
a collapsed interval imposes no condition. Define
\[
 f_\phi:\{1,\ldots,m\}\longrightarrow\{1,\ldots,n\},\qquad
 f_\phi(j)=\text{the unique }i\text{ with }\phi(i-1)<j\leq\phi(i).
\]
This is nondecreasing and sign-preserving. Conversely such a map $f$
defines
$\phi_f(i)=|\{j:f(j)\leq i\}|$. These two formulas are inverse.
They also cover the empty words: the map from a nonempty letter set to
an empty one does not exist, while the map from an empty letter set is
unique. If $\psi:I_\rho\to I_\sigma$ is another type map, the defining
interval inequalities give
$f_{\psi\phi}=f_\phi f_\psi$. Therefore these bijections form the
asserted isomorphism of categories with the opposite variance.

Let $g:q\to\sigma$ be any arrow in the opposite category. Since
adjacent letters of $q$ have opposite signs, sign preservation and
nondecreasingness imply
$g(1)<g(2)<g(3)$. Hence $ge$ is injective on the two letters of $p$,
whereas $hc$ is constant for every $h:q\to\sigma$. Thus $ge=hc$ is
impossible, proving failure of condition I.

For each $\sigma$ define an $R$-linear functional
\[
 \lambda_\sigma:R[\mathsf{Ty}^{\mathrm{op}}(p,\sigma)]\to R
\]
by giving a basis arrow value $1$ if it is injective and $0$ if it is
constant. These are the only possibilities for a map out of $p$.
For every $g:q\to\sigma$ we have
$\lambda_\sigma(ge-gc)=1$. Consequently for every finite linear
combination $t\in R[\mathsf{Ty}^{\mathrm{op}}(q,\sigma)]$,
\[
 \lambda_\sigma(t(e-c))=\operatorname{aug}(t).
\]
If $te=tc$, its augmentation is zero. It cannot also be $1$ unless
$R$ is the zero ring. This is an exact integral dual certificate,
valid for all targets and hence for all nonzero coefficient rings.

For the stronger homological claim, work over $\mathbb Z$ and put
$P_x=\mathbb Z[\mathsf{Ty}^{\mathrm{op}}(x,-)]$,
$d:P_q\to P_p$, $d(t)=t(e-c)$, $K=\ker d$, and $J=\operatorname{im}d$.
The objectwise free-group functor preserves the defining set
coequalizer, so $\operatorname{coker}d=\mathbb Z[H]$. Representables
are projective and their colimits are canonically $\mathbb Z$ by
augmentation. The functional identity shows that the map
$\operatorname{colim}K\to\operatorname{colim}P_q=\mathbb Z$ is zero,
even when its input is a finite sum of elements represented at different
objects. The exact sequence
$K\to P_q\to J\to0$ therefore gives
$\operatorname{colim}J\cong\mathbb Z$. The map
$\operatorname{colim}J\to\operatorname{colim}P_p$ is zero because the
augmentation of every $t(e-c)$ is zero. The derived-colimit exact
sequence of $0\to J\to P_p\to\mathbb Z[H]\to0$ now gives
\[
 L_1\operatorname{colim}\mathbb Z[H]
 \cong\operatorname{colim}J\cong\mathbb Z.
\]
The bar complex for the permutation diagram $\mathbb Z[H]$ is exactly
the integral nerve chain complex of its category of elements, proving
the claim. This explicitly witnesses failure of both totality and
exactness and does not infer an infinite conclusion from finite tests.

The program \texttt{code/finite\_ty\_verify.py} independently checks the
two descriptions on all $961$ ordered pairs of words of length at most
four, the contravariant composition bijection on $2731$ composable pairs
of type maps for words of length at most three, and the integral dual
certificate on all $7937$ arrows from $q$ to the $1023$ words of length
at most nine. These computations audit the formulas; the preceding
proof applies without a length bound.
\end{proof}

\subsection{Constraints on a minimal separator: Lead 9}

\begin{proposition}[Proved: Lead 9]\label{FT:minimal-family}
Suppose $A$ is a small total category that is not exact, and let
$a_0,\ldots,a_k:p\to q$ be a family without an integral affine equalizer.
Assume $k$ is least among all such families in $A$. Then:
\begin{enumerate}
\item $A(q,p)=\varnothing$.
\item The finite cardinalities $|A(p,x)|$ have no uniform finite bound;
in symbols, $\sup_x|A(p,x)|=\infty$ in the sense of unboundedness in the
natural numbers, allowing infinite values.
\item The full cosieve
$V_p=\{x\in A:A(p,x)\ne\varnothing\}$ is again a separator, with the
same exhibited family. Its object $p$ is weakly initial, and the minimal
integer $k$ remains the same after restriction.
\item For every arrow $g:q\to x$, the arrows
$ga_0,\ldots,ga_k:p\to x$ are pairwise distinct.
\end{enumerate}
\end{proposition}
\begin{proof}
Totality gives conditions I and IV. We use precisely the two exclusions
proved in the opened source \cite[Corollary~4.3 and Proposition~4.5]{IsbellMitchell1976}.
Corollary~4.3 first simultaneously equalizes a finite collection of
endomorphisms with the identity, by applying the fixed-point theorem to
endomorphisms of a representable act. If $b:q\to p$ existed, apply it to
$ba_0,\ldots,ba_k,1_p$. It gives an ordinary arrow $u:p\to r$ with
$uba_i=u$ for every $i$. The ordinary arrow $ub:q\to r$ would already
be an affine equalizer, proving (1).

For (2), the representable $A$-set $A(p,-)$ is connected because its
category of elements has the initial object $(p,1_p)$. If its values had
a uniform finite bound, Proposition~4.5 would imply condition II for this
act. That proposition forms the congruence identifying two elements
whenever an ordinary arrow equalizes them; condition I makes this a
congruence. All arrows act injectively on its quotient. At an object
where the finite cardinality is maximal, the maps to subsequent objects
are bijections, and the necklace form of IV forces the quotient to have
one element there. The conclusion is exactly ordinary equalizability
of every pair in the original representable. Apply it successively to
the finitely many images of $a_0,\ldots,a_k$; composing the resulting
ordinary arrows gives an ordinary simultaneous equalizer, a
contradiction. This verifies that the cited result applies to the one
representable involved here, rather than requiring a bound on every
hom-set of $A$.

For (3), $V_p$ is nonempty and is closed under targets. It is connected:
condition I applied to a span $p\to x$, $p\to y$ gives a common target
inside $V_p$. A covariant set-valued functor on $V_p$ extends to $A$ by
the empty set off $V_p$; arrows entering $V_p$ have the unique map from
an empty set, and no arrow leaves it. Its category of elements is
literally unchanged. Hence totality of $A$ proves totality of $V_p$.
Any affine equalizer of the displayed family in $V_p$ would also be one
in $A$, so $V_p$ is not exact. Conversely, if a family of arrows in
$V_p$ has an affine equalizer in $A$, every target occurring in its
finite support is in $V_p$, because it receives an arrow from its source
in the cosieve. Thus it is already an affine equalizer in $V_p$.
All smaller families therefore remain equalizable. Every object of
$V_p$ receives a map from $p$, which is exactly weak initiality.
This is the full proof of the cosieve step used from TOTEX D.23.

For (4), suppose $ga_i=ga_j$ for some $i\ne j$. The displayed family at
$x$ then has at most $k$ distinct arrows. By the global minimality of
$k$ there is an affine $t:x\to y$ equalizing all its distinct members.
The composite $tg\in\mathbb Z[A(q,y)]$ has augmentation one, because
$g$ is an ordinary arrow, and satisfies $(tg)a_i=(tg)a_0$ for every
$i$. This contradicts the choice of the original family. This argument
uses minimality among all objects and families of $A$, as stipulated in
the lead; minimality merely within a fixed target would not suffice.
\end{proof}
